L'intégration d'outils d'intelligence artificielle dans les bibliothèques d'orchestre et leurs plateformes numériques pour la transformation de leurs ressources en données exploitables doit, comme nous l'avons vu, être encadrée. 
Dans la continuité de la politique documentaire, étudiée au chapitre 4, et des enjeux de normalisation des données, étudiés au chapitre 5, le recours à l'IA suppose en effet un cadre rigoureux. 

Dans le domaine des plateformes numériques musicales, comme pour celui des industries culturelles et créatives en général, l'intégration de l'IA se heurte à des contraintes juridiques et éthiques fondamentales quant à l'intégrité du patrimoine, au respect du droit d'auteur ou encore à la protection des personnes et de leur vie privée. 


\subsection*{Le cadre juridique des ressources patrimoniales musicales à l'ère de l'intelligence artificielle}

L'adoption d'outils d'IA dans les plateformes numériques comme LaDocumenta.eu repose sur des volumes importants de données. Dans le domaine de la pratique historiquement informée de la musique, cette nouvelle donne implique de réinterroger le statut de la ressource patrimoniale musicale face aux outils d'IA générative, textuelle ou audio.

\subsubsection*{Droit d'auteur et musique ancienne : la facilité trompeuse du domaine public}

Sans chercher une exhaustivité juridique qui n'aurait pas sa place dans ce mémoire, il est toutefois fondamental de poser ici les jalons des limites juridiques émanant des utilisations numériques des ressources patrimoniales musicales des bibliothèques musicales, et de leur interface numérique.
Si les oeuvres de la musique historiquement informée des périodes Renaissance, baroque, classique et romantique, appartiennent au domaine public\footnote{Code de la propriété intellectuelle (CPI), Art. L. 123-1 : \enquote{Article L123-1 Modifié par Loi n°97-283 du 27 mars 1997 - art. 5 () JORF 28 mars 1997 en vigueur le 1er juillet 1995
		\textit{L'auteur jouit, sa vie durant, du droit exclusif d'exploiter son oeuvre sous quelque forme que ce soit et d'en tirer un profit pécuniaire.}
		\textit{Au décès de l'auteur, ce droit persiste au bénéfice de ses ayants droit pendant l'année civile en cours et les soixante-dix années qui suivent.}}
}, les éditions qui en émanent ne sont pas nécessairement dans le domaine public. 

En effet, une ressource musicale peut dépendre de trois états juridiques distincts : 

\begin{itemize}
	
	\item\textbf{L'oeuvre originale et le texte musical} : dans ce cas, la composition n'est plus couverte par le droit d'auteur à partir de soixante-dix ans après la mort du compositeur. Une précision importante est à ajouter concernant le texte de l'oeuvre. En effet, soixante-dix ans après la mort du compositeur, le matériel musical (partition, etc) entre dans le domaine public aussi. 

	\item \textbf{L'édition critique (\textit{Urtext}) ou moderne de l'œuvre}~: distincts de l'œuvre originale, les choix éditoriaux, la reconstitution de parties manquantes ou l'apparat critique établis par un musicologue ou une maison d'édition peuvent bénéficier d'une protection par le droit d'auteur, à condition qu'ils témoignent d'un apport créatif réel. L'édition est alors protégée pendant soixante-dix ans après le décès de l'éditeur scientifique. Dans le cas d'un simple travail d'érudition ou d'établissement scientifique du texte sans création originale, la protection peut être restreinte à un droit voisin de vingt-cinq ans à compter de la publication. Ce cas de figure s'est concrètement présenté pour l'ensemble Insula Orchestra lors de la programmation d'une symphonie de Louise Farrenc : bien que la composition soit dans le domaine public, le matériel d'orchestre récent demeurait sous droits d'auteur en raison du travail d'édition contemporain. L'orchestre se trouvait dès lors contraint de louer ce matériel protégé.
	
	\item \textbf{L'interprétation de l'oeuvre en concert, et ses éventuelles captations audio} : cette configuration est encadrée quant à elle par des droits dits voisins, c'est-à-dire des droits \enquote{accordés à différentes catégories de personnes qui gravitent autour des auteurs. Il s'agit des artistes-interprètes, des producteurs de phonogrammes et de vidéogrammes, des entreprises de communication audiovisuelle et de télédiffusion par satellite et de retransmission par câble\footnote{Source : Rubrique des droits voisins de Dalloz.}}. Les enregistrements sonores sont également protégés pendant soixante-dix ans.
\end{itemize}

À ces considérations générales qui régissent le cadre juridique des œuvres musicales s'ajoute une disposition spécifique issue du droit européen et transposée en droit national : l'exception de fouille de textes et de données (Text and Data Mining ou TDM). Au cœur des enjeux contemporains liés à la constitution de corpus d'entraînement, ce cadre légal est issu à l’origine de la directive européenne (UE) 2019/790 et s'applique aujourd'hui en France par l'article L. 122-5-3 du Code de la propriété intellectuelle
\footnote
{Article L. 122-5-3 du Code de la propriété intellectuelle :
	\enquote{
		\textit{Des copies ou reproductions numériques d'œuvres auxquelles il a été accédé de manière licite peuvent être réalisées sans autorisation des auteurs en vue de fouilles de textes et de données menées à bien aux seules fins de la recherche scientifique par les organismes de recherche, les bibliothèques accessibles au public, les musées, les services d'archives ou les institutions dépositaires du patrimoine cinématographique, audiovisuel ou sonore, ou pour leur compte et à leur demande par d'autres personnes, y compris dans le cadre d'un partenariat sans but lucratif avec des acteurs privés}.
	}
}.

Cette exception correspond à un axe de recherche scientifique. Il permet aux organismes de recherche et aux institutions du patrimoine culturel de faire des fouilles automatisées sur des contenus numériques et ce de manière légale. L'accès aux données et leur traitement est alors de plein-droit, sans autorisation préalable nécessaire de l'auteur, ni compensation financière.


\subsubsection*{La propriété intellectuelle à l'ère de l'IA}

Par ailleurs, au-delà des œuvres musicales et de leurs éditions que conserve la plateforme, se pose la question de la propriété intellectuelle des contenus produits avec l'aide de l'IA, comme par exemple les métadonnées enrichies.

Ces questionnements occupent de plus en plus les débats juridiques, notamment au Parlement européen où la résolution \enquote{Le droit d'auteur et l'intelligence artificielle générative – opportunités et défis}\footnote{\textsc{Parlement européen}, \emph{Résolution du 10 mars 2026 sur le droit d’auteur et l’intelligence artificielle générative – opportunités et défis} (2025/2058(INI)), texte adopté P10\_TA(2026)0066.} a été adoptée, même si, aujourd'hui, ces métadonnées ne bénéficient pas de la protection par le droit d’auteur en raison du critère de la \enquote{prépondérance du créateur humain}\footnote{\textsc{Ministère de l'Économie, des Finances et de la Souveraineté industrielle et numérique} — Agence du patrimoine matériel de l'État (APIE), \enquote{Quel droit d'auteur à l'ère de l'intelligence artificielle générative ?}, \emph{Portail de l'Économie et des Finances}, [en ligne] :~\url{https://www.economie.gouv.fr/apie/quel-droit-dauteur-lere-de-lintelligence-artificielle-generative} (consulté le 28 août 2026).}.



Par ailleurs, dans l'esprit de la science ouverte, étudiée au chapitre 1, les données produites ou enrichies par l’IA  au sein du réseau doivent être diffusées sous des licences ouvertes, garantissant leur libre réutilisation par les chercheurs et musiciens du monde entier. La licence Creative Commons BY, par exemple, permet ainsi de réutiliser et modifier du contenu à la condition de citer la source originale. 



\subsection*{Les enjeux éthiques des bibliothèques musicales à l'ère de l'IA : repenser la recherche et les métiers des bibliothèques.}

Les enjeux éthiques, abordés dans le cadre de l'émergence de l'IA dans les domaines des industries culturelles et créatives, sont complémentaires au cadre juridique que nous venons d'étudier. 

La réflexion éthique sur l'intelligence artificielle garantit que l'IA reste un outil d'assistance, respectueux des missions des bibliothèques et de leurs usagers. 

Afin d'aborder les spécificités de l'éthique - à l'ère de l'IA - en bibliothèques, notre réflexion s'est appuyée sur les principes de la BnF et des instances référentes, européenne et mondiale, c'est-à-dire respectivement la CENL (\textit{The Conference of European National Librarians}) et la communauté IA4LAM (\textit{Artificial Intelligence for Library, Archives, Museums}). 

Dans un document collaboratif portant sur les enjeux de l'intelligence articielle, la Bibliothèque nationale de France présente sa démarche éthique rigoureuse guidée par cinq axes\footnote{Ces axes sont détaillés dans un document de la stratégie d'intelligence articielle de la BnF, en ligne à cette adresse : \url{https://www.bnf.fr/fr/ethique-et-ia-la-bnf}} :

\begin{itemize}
	\item \enquote{recourir à l'IA au service de l'intérêt général}
	\item \enquote{contribuer au renforcement d'une IA de confiance}
	\item \enquote{veiller au respect des droits des personnes}
	\item \enquote{lutter contre les biais et discriminations reproduits par l'IA}
	\item \enquote{réduire l'impact environnemental des IA}
\end{itemize}



Les enjeux éthiques des bibliothèques numériques musicales reprennent les mêmes principes que ceux préconisés par la BnF. 
En bibliothèques musicales numériques, l’IA doit également être au service de l’intérêt général, comme dans toute bibliothèque. Appliqués aux infrastructures spécialisées comme LaDocumenta.eu, ces principes sont donc au service de la recherche en musicologie (HIP), des musiciens et du public. Le principe de complémentarité est celui qui décrit le mieux les dispositions éthiques de ces structures. D'après ce principe, l’IA doit intervenir comme un assistant pour les tâches répétitives (comme la saisie des métadonnées) qui permet aux bibliothécaire de gagner du temps précieux. L’IA peut également être un levier d’aide à la décision dans le cadre de la politique documentaire (chapitre 4), pour la sélection et l’enrichissement des fonds, avec une validation humaine obligatoire ensuite.



Par ailleurs, le rôle des bibliothèques est de lutter contre les biais. En bibliothèques musicales, cela peut se traduire par une vérification humaine de la diversité des contenus proposés par la plateforme, notamment en évitant la surreprésentation de chefs d’oeuvre grand public (les \enquote{pépites} dans la plateforme LaDocumenta.eu) aux dépens d’oeuvres de valeur n’ayant pas connu le même succès. La redécouverte des œuvres anciennes oubliées étant un des principes de la pratique historiquement informée, la lutte contre ce biais est très importante pour les chercheurs et musiciens.



En outre, les enjeux de transparence et de traçabilité de la donnée, garantissant le développement d’une IA de confiance, sont des domaines où les bibliothèques doivent se montrer pionnières. Dans un souci d’honnêteté intellectuelle, les bibliothèques ont un rôle d’exemple à tenir en documentant toute intervention de l’IA, dans l’enrichissement automatique des métadonnées par exemple, ou dans un traitement de ressource musicale en OMR par exemple. Chercheurs, musiciens et public sont censés pouvoir distinguer la contribution humaine d’une tâche réalisée par l’IA. Ces enjeux se placent au centre des débats actuels de transformation de la recherche universitaire, évoqués au chapitre 7\footnote{Cases (Anne-Sophie ), Fournier (Christophe), « Défis éthiques des IA génératives pour l’enseignement supérieur », Revue Management et Data Science, 2025. En ligne : \url{https://management-datascience.org/articles/36833/}}.



Enfin, comme dans toute organisation, régie par la Responsabilité Sociétale des Organisations (RSO), l’éthique de l’utilisation de l’IA intègre des enjeux de protection des personnes et de l’environnement dans ses principes. Ces enjeux, de protection des données personnelles, de souveraineté culturelle et de sobriété numérique, tels que décrits par la BnF ont pour objectif de construire une IA responsable.
Dans le processus de transformation des ressources musicales en données exploitables, le rôle de la plateforme numérique offre des solutions originales et prometteuses. Toutefois les enjeux éthiques de protection des données des usagers doivent être conformes aux pratiques réglementées du RGPD (Réglement Général de Protection des données). Assurer la protection des usagers, contributeur comme public, est précisement le coeur de la mission de protection de la personne énoncée par les principes éthiques de la BnF.



Dès lors, si la protection des personnes est un enjeu important des institutions culturelles utilisant l’IA, la souveraineté culturelle et la protection des données l’est tout autant. Les institutions doivent dès lors promouvoir les modèles souverains et ouverts, ancrés dans les réseaux de recherche européens. Ces solutions prometteuses offertes par les plateformes numériques ne sont pas exemptes de l’impact énergétique considérable du traitement des masses de données. Il relève de la responsabilité des bibliothèques numériques musicales d’évaluer l’empreinte carbone de leurs projets IA et de tenir compte de cet élément pour le choix de ceux-ci, comme cela est préconisé dans le texte de référence de la BnF.



Ainsi, en articulant les aspects juridiques et éthiques de l’intelligence artificielle, les plateformes numériques musicales telles que LaDocumenta.eu affirment le rôle du bibliothécaire musical comme un garant des bonnes pratiques éthiques et juridiques des bibliothèques. A l’ère de l’IA, la valeur des bibliothèques musicales, et notamment des consortiums qui les rassemblent, se mesure dans leur capacité à proposer un environnement de données structurées, interopérables et exploitables, souverain, au service de la recherche musicologique et de ses réseaux mais aussi de l’interprétation de pratique historiquement informée et de son public.
